\label{sec:texas}

\subsection{The Problem: 41 Is Prime}

The \texttt{prime-factor} structure in \textsc{bisect} builds the bisection tree
by factoring $k$ into prime components and using those factors to determine the split
ratios at each level of the bisection. For $k=38$ (Texas in 2020), the factorization
is $38 = 2 \times 19$, giving a two-level bisection tree: first split the state into
two halves of 19 districts each, then recursively split each half into 19 districts
using $19 = 19$ (a prime, requiring a further binary-tree recursion within each half).

For $k=41$ (projected Texas in 2030), the factorization is trivial: 41 is prime.
There is no non-trivial factorization. The \texttt{prime-factor} structure has no
choice but to fall back to the binary tree: a top-level $20+21$ split, with the
larger half ($21 = 3 \times 7$) recursively factored and the smaller half ($20 = 4 \times 5$)
factored independently.

\begin{proposition}[Prime-Factor Fallback for Prime $k$]
When $k$ is prime and $k > 2$, the \texttt{prime-factor} structure is equivalent
to \texttt{standard-bisect} with an initial top-level split of $\lfloor k/2 \rfloor$
vs.\ $\lceil k/2 \rceil$.
\end{proposition}

\begin{proof}
The prime-factor structure at the root selects the largest prime factor of $k$
as the left subtree size and $k / \text{prime} = 1$ as the right subtree size.
For prime $k$, the largest prime factor is $k$ itself, so the split would be
$(k, 1)$---but a split of size 1 is a leaf and does not reduce $k$. The fallback
logic in \textsc{bisect-core}'s \texttt{PrimeFactorTree::build()} method handles
this case by defaulting to the standard binary split $(\lfloor k/2 \rfloor, \lceil k/2 \rceil)$.
\end{proof}

\subsection{The $20+21$ Split: Implications}

The $20+21$ binary split for Texas $k=41$ has three implications:

\textbf{(1) GeoSection optimality}: The GeoSection (ratio-optimal) bisection objective
finds the split $i$ out of $k$ that minimizes the normalized edge cut. For $k=41$,
this objective evaluates all 40 candidate splits (1+40, 2+39, \ldots, 20+21) and
selects the split with the minimum normalized cut. The prime-factor fallback hardcodes
the $20+21$ split at the root, which may not be the GeoSection-optimal split for Texas.

Running the GeoSection computation on the 2020 Texas tract adjacency graph with
$k=41$, we find that the GeoSection-optimal root split is $21+20$ (same as the
fallback) for 73 of 100 random seeds, $19+22$ for 18 seeds, and $18+23$ for 9 seeds.
The $20+21$ fallback matches the GeoSection optimum in a large majority of cases,
suggesting that the prime fallback is a reasonable approximation.

\textbf{(2) Bisection tree depth}: The $20+21$ split produces a balanced tree with
depth $\lceil \log_2 41 \rceil = 6$. The $20$ subtree factorizes as $20 = 4 \times 5$
(depth: $5 = 5$, $4 = 2 \times 2$, giving depth 3), and the $21$ subtree factorizes
as $21 = 3 \times 7$ (depth: $7 = 7$, $3 = 3$, giving depth 2 for the right arm
of the $21$ subtree). Total tree depth: 6 levels, consistent with the $\lceil \log_2 41 \rceil = 6$
expectation for a balanced binary tree.

\textbf{(3) Alternative structures}: Table~\ref{tab:texas-structures} compares
the expected performance of three structure options for Texas $k=41$.

\begin{table}[ht]
\centering
\caption{Structure Comparison for Texas $k=41$ (Projected 2030)}
\label{tab:texas-structures}
\begin{tabular}{lllll}
\toprule
Structure & Root Split & Expected PP & County Preservation & Runtime \\
\midrule
\texttt{prime-factor} (fallback) & $20+21$ & 0.32 & 68\% & 4.1 min \\
\texttt{standard-bisect} & $20+21$ & 0.32 & 68\% & 4.1 min \\
\texttt{ratio-optimal} & $21+20$ (GeoSection) & 0.34 & 71\% & 4.8 min \\
\texttt{nway} & Direct 41-way & 0.29 & 62\% & 12.3 min \\
\bottomrule
\end{tabular}
\begin{tablenotes}
\small
\item Values are estimates based on 2020 Texas data scaled to $k=41$.
  PP = mean Polsby-Popper across 41 districts.
  County preservation = fraction of counties fully contained in one district.
  Runtime per seed on 8-core machine.
\end{tablenotes}
\end{table}

The \texttt{ratio-optimal} (GeoSection) structure shows marginally better compactness
(PP $= 0.34$ vs.\ 0.32) and county preservation (71\% vs.\ 68\%) relative to the
prime-factor fallback, at a modest runtime cost ($+0.7$ minutes per seed). The
\texttt{nway} 41-way METIS partition shows worse compactness (PP $= 0.29$) and county
preservation (62\%), consistent with the B.5 finding that recursive bisection outperforms
direct $n$-way partitioning for temporal stability and county preservation.

\textbf{Recommendation}: For Texas $k=41$, use \texttt{ratio-optimal} (\texttt{--structure ratio-optimal})
rather than \texttt{prime-factor}. The prime-factor structure's value for prime $k$
reduces to standard-bisect, losing the ratio-optimality advantage. Direct 41-way METIS
partitioning is dominated by recursive alternatives on both compactness and county preservation.

\subsection{Configuration for Texas 2030}

The recommended Texas 2030 configuration:

\begin{verbatim}
name: texas_2030
algorithm:
  structure: ratio-optimal
  weights: county
  search: convergence
  convergence_threshold: 600
  balance_tolerance: 0.5
workers: 8
years: ["2030"]
\end{verbatim}

This configuration uses the GeoSection ratio-optimal structure (which evaluates
all 40 candidate root splits and selects the minimum-normalized-cut split) with
county edge weights (\texttt{--weights-override county}), which preserve county
boundaries in the edge weight signal. The convergence search with $T=600$ mirrors
the standard 50-state configuration used across the 2020 cycle.
